\section{能量、亲密关系与记忆程序}

上一章把组织理解为共同大脑和共同心脏，本章进入更底层的能量来源。本节先看访谈里最不像技术、但最能解释李想方法论的一部分：能量来自成长、来自亲密关系、来自被需要和需要别人。他说自己需要孩子、爱人和一级部门负责人，甚至他们对自己的重要性超过自己对他们的重要性。这种表达不是情感鸡汤，而是把关系看成一种反馈系统：亲密关系让人看到自己的不足，也让人持续变好。

\begin{figure}[H]
\centering
\includegraphics[width=0.96\textwidth]{energy-relationship-loop.png}
\caption{能量与亲密关系：关系不是消耗项，而是让人变得更好的反馈系统。自制概念图，依据 02:09:26--02:28:20 对谈内容整理。}
\end{figure}

\begin{knowledgebox}{读图：能量不是兴奋感，而是持续反馈}
需求、连接、反馈、能量和成长形成循环。亲密关系之所以重要，是因为它能提供真实反馈；真实反馈会暴露不足；接受不足并变好，才会产生持续能量。
\end{knowledgebox}

\subsection{亲密关系的边界}

前面强调亲密关系会带来能量，本节补上同样重要的边界条件。李想也强调，不要构建太多亲密关系。亲密关系不是泛泛社交，而是直系亲属、少数多年朋友、以及工作中一起扛责任的人。能带来伤害的，往往也只有亲密关系，因为它真正进入了人的需求和脆弱处。这个观点可以迁移到组织：高质量关系有能量，低质量关系会消耗能量。

\begin{importantbox}{亲密关系的定义}
亲密关系不是聊天频率高，也不是社交关系近，而是彼此能真实需要、真实反馈、共同承担后果，并在长期中让对方变好。
\end{importantbox}

\subsection{记忆程序：如果可以修改一段人生代码}

访谈里还出现了“记忆程序”的问题：如果能删除或修改记忆，会选择什么。李想的回答并不把痛苦记忆简单视为坏事，因为很多记忆塑造了今天的方法论、关系和判断。这里可以把“记忆”理解为人的训练数据：它不只是存档，也是行为策略和价值判断的来源。

\begin{knowledgebox}{记忆作为人的训练数据}
AI 模型通过数据形成参数，人则通过记忆、关系和行动形成方法论。删除痛苦记忆看似降低损耗，但也可能删除能力来源。重要的是把记忆消化成方法，而不是被记忆反复触发。
\end{knowledgebox}

\subsection{本章小结}

能量和亲密关系不是访谈里的“人文插曲”，而是李想理解组织和智慧的底层变量。没有能量，方法论无法长期执行；没有真实关系，人很难看见自己的盲区。

